% 本文件由 LaTeX 排版 Agent 根据有效 translation.json 自动生成。
% author=Cyprian of Carthage; work_id=0508; title=论异端者的圣洗

\chapter{论异端者的圣洗}

% structure: p0001 source=h1 target=omitted reason=page-title-already-rendered-by-parent

\Emph{\Place{迦太基}第七次公会议，由\Person{西彼廉}主持。}\par

\Emph{论异端者的圣洗。八十七位主教关于异端者圣洗的判决。}\par

序言。——当罗马主教\Person{斯德望}以书信谴责非洲公会议关于异端者圣洗的法令时，\Person{西彼廉}毫不迟延，在\Place{迦太基}召开了另一次人数更多的公会议。他因此从\Place{阿非利加}、\Place{努米底亚}和\Place{毛里塔尼亚}召集了八十七位主教，他们于公元258年九月朔日在\Place{迦太基}集会，于是举行了这第三次关于同一圣洗问题的公会议；会议伊始，在双方的书信被宣读之后，\Person{西彼廉}含蓄地谴责了\Person{斯德望}的僭越。\par

\Person{西彼廉}说：你们已经听到了，我亲爱的同僚们，我们的同主教\Person{尤巴伊安努斯}写信给我，就异端者非法和亵渎的圣洗咨询我这贫乏的智慧，以及我写给他的回信，我等过去、现在和经常都这样决定，即凡来到教会的异端者必须由教会的圣洗予以圣洗和圣化。此外，\Person{尤巴伊安努斯}的另一封书信也已向你们宣读，其中他以真诚和虔诚的热忱回复我的信，不仅同意我所说的，而且承认从中得到教导，并为此表示感谢。剩下的，就是关于同一件事，我们每个人都要提出自己的想法，不要判断任何人，也不要因与他人意见不同而拒绝其共融的权利。因为我们中没有人自称是主教的主教，也没有人以暴虐的恐怖强迫其同僚必须服从；因为每一位主教，根据其自由和权力的许可，都有自己适当的判断权，既不能受他人判断，也不能判断他人。但我们都等待我们主\Person{耶稣基督}的审判，惟有祂有权柄在治理祂的教会时赏赐我们，并在我们的行为中审判我们。\par

\Place{比尔塔}的\Person{凯奇利乌斯}说：我知道教会内只有一个圣洗，教会外没有。这一个圣洗将在这里，这里有真实的希望和确定的信德。因为经上记载：\BibleQuote{一个信德，一个希望，一个圣洗；}(\BibleRef{弗 4:5})不在异端者那里，那里没有希望，信德是虚假的，所有事都是靠着谎言进行的；那里是魔鬼在驱魔；那里是口出毒疮的人提出圣事的质询；无信者赋予信德；邪恶者赐予罪的赦免；敌基督以基督之名施洗；被天主诅咒的人祝福；死人应许生命；不平安的人赐予平安；亵渎者呼求天主；亵圣者行使司祭职；亵渎神明的人设立祭台。除了这一切，还有这恶事：魔鬼的司祭竟敢举行圣体圣事；否则，就让那些支持他们的人说，所有这些关于异端者的事都是假的。请看，教会被迫同意什么样的事，被迫在没有圣洗、没有罪的赦免的情况下与他们共融。弟兄们，我们应当逃避和避免这事，远离如此大的邪恶，持守那唯独由主赐予教会的唯一圣洗。\par

\Place{米斯吉尔帕}的\Person{普里穆斯}说：我决定，凡从异端来到我们这里的人，都必须受圣洗。因为他若以为在那里已受了圣洗，便是徒然，因为只有在教会内才有那唯一真实的圣洗；因为天主不仅是一位，信德也是一，教会也是一，其中存在着唯一的圣洗、圣洁及其他一切。因为凡在教会之外所做的，都没有救恩的实效。\par

\Place{阿德鲁梅图姆}的\Person{波利卡普}说：那些认可异端者圣洗的人，使我们的圣洗归于无效。\par

\Place{塔穆加达}的\Person{诺瓦图斯}说：虽然我们知道所有圣经都为救恩的圣洗作证，但我们仍应宣认我们的信德，即异端者和分裂者来到教会，他们似乎受过虚假的圣洗，应当在水泉中受圣洗；因此，按照圣经的见证，并按照我们同僚、至圣记忆者的谕令，所有归向教会的分裂者和异端者必须受圣洗；而且，那些似乎已受祝圣的人，应当被接纳为平信徒。\par

\Person{Nemesianus of Thubunae}说：异端者和分裂者所施行的洗礼不是真正的洗礼，这在圣经中到处都有宣告，因为他们的首领本身就是假基督和假先知，正如主通过\Person{撒罗满}所说：\Quote{那信赖虚假事物的人，是在喂养风；并且同样，他跟随飞鸟的路径。因为他离弃了自己葡萄园的道路，从自己小田地的路径中迷失了。但他行经无路、干旱和注定干渴之地；此外，他在手中收集无果之物。}又说：\BibleQuote{你要远离陌生的水，不要从别人的泉源中喝，这样你才能长寿；并且生命的年日也要加给你。} \BibleRef{箴 9:19} 而在福音中，我们的主耶稣基督用祂神圣的声音说：\BibleQuote{人若不从水和圣神重生，不能进入天主的国。} \BibleRef{若 3:5} 这就是那从起初运行在水面上的圣神；因为圣神不能没有水而运作，水也不能没有圣神而运作。因此，有些人错误地为自己解释，说藉着覆手他们领受了圣神，就这样被接纳了，但事实上他们显然应当在天主教会中藉着两件圣事重生。然后他们才能成为天主的子女，正如宗徒所说：\BibleQuote{以和平的连系，竭力保持圣神所赐的合一。只有一个身体和一个圣神，正如你们蒙召同有一个希望一样；一主、一信、一洗、一天主。} \BibleRef{弗 4:3} 这一切都是天主教会所说的。再者，主在福音中说：\BibleQuote{由肉生的属于肉，由圣神生的属于神；因为天主是神，并且由天主而生。} \BibleRef{若 3:6} 因此，所有异端者和分裂者所行的一切都是属血肉的，正如宗徒所说：\BibleQuote{本性私欲的作为是显而易见的：即淫乱、不洁、放荡、崇拜偶像、行邪术、仇恨、竞争、嫉妒、愤怒、纷争、异端、嫉妒、醉酒、宴乐等类。我以前说过，现在又预先说：行这样事的人，决不能承受天主的国。} \BibleRef{迦 5:19} 因此，宗徒与一切恶人一同谴责那些制造分裂的人，即分裂者和异端者。除非他们在天主教会——那唯一的教会——中领受救恩的洗礼，否则不能得救，反而要在主基督的审判中与属血肉的人一同被定罪。\par

\Person{Januarius of Lambesis}说：根据圣经的权威，我判定所有异端者必须受洗，并这样被接纳进入圣教会。\par

\Person{Lucius of Castra Galbae}说：既然主在福音中说：\BibleQuote{你们是地上的盐；盐若失了味，可用什么使它再咸呢？它便毫无用途，只好抛在外面，任人践踏。} \BibleRef{玛 5:13} 再者，在祂复活后，派遣祂的宗徒时，祂吩咐他们说：\BibleQuote{天上地下的一切权柄都交给了我。你们要去使万民成为门徒，因父、及子、及圣神之名给他们授洗。} \BibleRef{玛 28:18} 因此，既然异端者——即基督的敌人——显然没有对圣事的正确宣认；况且分裂者不能用属神的智慧调味他人，因为他们自己因离开了那唯一的教会，失去了味，成为与之对立——那么，就应按所写的去做：\Quote{违法者的房屋需要洁净。} 结果是，那些由与教会对立的人施洗而受玷污的人，必须先被洁净，然后才受洗。\par

\Person{Crescens of Cirta}说：在如此众多至圣的同僚司铎的集会中，既然我们至爱的西彼廉写给尤巴亚努斯以及也给斯德望的书信已被宣读，其中包含着如此多来自神圣圣经的神圣见证，以致我们有理由藉着天主的恩宠，同心合意地同意它们。我判定，所有愿意来到天主教会中的异端者和分裂者，除非他们先被驱魔和受洗，不得进入；但那些先前已在天主教会中受洗的人除外，这些人则应以覆手的方式与教会的补赎和好。\par

\Person{Nicomedes of Segermae}说：我的意见是，来到教会的异端者应该受洗，因为他们在罪人中无法获得罪的赦免。\par

\Person{Munnulus of Girba}说：我们母亲天主教会会的真理，弟兄们，常与我们同在，现在仍与我们同在，尤其是在洗礼的天主圣三中，正如主所说：\BibleQuote{你们去，给万民授洗，因父、及子、及圣神之名。} \BibleRef{玛 28:19} 既然我们明确知道异端者既没有父，也没有子，也没有圣神，那么他们来到我们母亲天主教会时，应当真正重生并受洗；好使他们所患的毒瘤、诅咒的愤怒和错误的邪术，能被圣洁而属天的洗礼所圣化。\par

\Person{Secundinus of Cedias}说：既然我们的主基督说：\BibleQuote{不随同我的，就是反对我；} \BibleRef{玛 12:30} 而若望宗徒称那些离开教会的人为假基督——无疑是基督的敌人——那么被称为假基督的人不能施行救恩洗礼的恩宠。因此，我认为那些从异端者的罗网逃到教会的人，应由我们——因天主的屈尊就卑被称为天主朋友的人——来施洗。\par

\Person{Felix of Bagai}说：正如瞎子领瞎子，两人都掉在坑里；同样，异端者给异端者施洗，两人都同归于死亡。因此，异端者必须受洗而重生，免得我们这些活人与死人交往。\par

\Person{Polianus of Mileum}说：异端者在圣教会中受洗是正确的。\par

提奥革内斯（\Person{Theogenes}）从\Place{Hippo Regius}说：根据我们所领受的天主属天恩宠的圣事，我们相信在圣教会中只有一个圣洗。\par

达提乌斯（\Person{Dativus}）从\Place{Badis}说：按我们所能，我们不与异端者共融，除非他们在教会中受了圣洗并获得了罪的赦免。\par

苏克苏斯（\Person{Successus}）从\Place{Abbir Germaniciana}说：异端者要么什么都不能做，要么就能做一切。如果他们能施洗，他们也能赐予圣神。但如果他们不能赐予圣神，因为他们没有圣神，那么他们也不能属灵地施洗。因此我们判定异端者必须受洗。\par

福图纳图斯（\Person{Fortunatus}）从\Place{Tuccaboris}说：我们的主天主耶稣基督、圣父天主子与创造者，将祂的教会建立在磐石上，而非建立在异端上；并将施洗的权柄赐给了主教，而非异端者。因此，那些在教会之外并与基督对立、分散祂的羊群的人，不能施洗，因为他们在外面。\par

塞达图斯（\Person{Sedatus}）从\Place{Tuburbo}说：正如在教会中由司铎的祈祷所祝圣的水能洗去罪恶，同样，若被异端言论如同癌症般感染，它便堆积罪恶。因此我们必须以一切和平的力量努力，使任何被异端错误感染和玷污的人，不可拒绝接受教会唯一而真实的圣洗，凡未受此洗者，将成为天国的局外人。普里瓦提亚努斯（\Person{Privatianus}）从\Place{Sufetula}说：凡说异端者有施洗权柄的人，先要说谁创立了异端。因为如果异端出于天主，它也可能拥有天主的宽恕。但如果它不出于天主，它怎能拥有天主的恩宠，或将恩宠赐予任何人呢？\par

普里瓦图斯（\Person{Privatus}）从\Place{Sufes}说：凡赞同异端者圣洗的人，除了与异端者共融外，还做了什么？\par

荷尔滕西安努斯（\Person{Hortensianus}）从\Place{Lares}说：让这些傲慢的人，或那些偏袒异端者的人，思考有多少圣洗。我们为教会主张一个圣洗，这圣洗我们只知道在教会内。或者，他们怎能以基督的名义给任何人施洗，而基督自己宣称这些人是祂的敌人？\par

卡西乌斯（\Person{Cassius}）从\Place{Macomadae}说：既然不能有两个圣洗，凡将圣洗让给异端者的人，就是从他自身夺走了圣洗。因此我判定，可悲而败坏的异端者当他们开始来到教会时，必须受洗；并且当他们被神圣而属天的洗涤所洗，被生命之光照亮时，他们可以被接纳进入教会，不是作为敌人，而是作为被和好的人；不是作为外人，而是作为主信仰的家人；不是作为奸淫之子，而是作为天主的子女；不是出于错误，而是出于救恩；唯独那些一度忠信却被颠覆、从教会转向异端黑暗的人，必须借着按手而复原。\par

另一位雅努阿里乌斯（\Person{Januarius}）从\Place{Vicus Caesaris}说：如果错误不顺从真理，那么真理更不会同意错误；因此我们持守我们在其中主持的教会，以便为她独自主张她的圣洗，我们应当给那些教会未曾施洗的人施洗。\par

另一位塞昆迪努斯（\Person{Secundinus}）从\Place{Carpi}说：异端者是不是基督徒？如果是，他们为什么不在天主的教会内？如果不是，他们怎能使人成为基督徒？或者，主的言论将指向何处，当祂说：\Quote{不与我同在的就是反对我，不与我收集的就是分散。} \BibleRef{玛 12:30}由此清楚显明，对于陌生的儿女和敌基督的后裔，圣神不能仅借按手而降临，因为显然异端者没有圣洗。\par

维克多利库斯（\Person{Victoricus}）从\Place{Thabraca}说：如果异端者被允许施洗并赦罪，那么我们为什么要给他们打上恶名并称他们为异端者？\par

另一位费利克斯（\Person{Felix}）从\Place{Uthina}说：最神圣的同司铎们，没有人怀疑人的僭越无法比得上我们主耶稣基督可敬而庄严的威严。因此，记住危险，我们不仅应当遵守这一点，而且还应当借我们所有人的声音来确认它：凡来到慈母教会怀抱的异端者都应受洗，以便被长期败坏所玷污的异端心思，借着洗礼盆的圣化而得以净化，从而向善更新。\par

奎厄图斯（\Person{Quietus}）从\Place{Baruch}说：我们这些因信德而生活的人，应当以谨慎的遵守来服从那些预先为我们教诲所说的话。因为在\BibleBook{}{德训篇}中写着：\Quote{那从死者中受洗（又再触摸死者）的，他的洗涤有何益处呢？} \BibleRef{德 34:25}这无疑是指那些被异端者清洗的人，以及那些清洗他们的人。因为如果那些在他们中受洗的人借着罪的赦免获得了永生，他们为什么来教会？但如果从死者那里得不到救恩，并且因此，他们承认先前的错误，带着补赎回到真理，他们应当被唯一生命的圣洗所圣化，这圣洗是在天主教会中的。\par

卡斯图斯（\Person{Castus}）从\Place{Sicca}说：凡藐视真理而妄图跟随习俗的人，要么是对他的弟兄们（真理已向他们启示）怀有嫉妒和恶意，要么是对天主（教会借祂的灵感而得教导）忘恩负义。\par

欧克拉提乌斯（\Person{Euchratius}）从\Place{Thenae}说：天主和我们的主耶稣基督，亲口教导宗徒们，完全完成了我们的信仰、圣洗的恩宠以及教会法的规则，说：\Quote{你们去教导万民，因父及子及圣神之名给他们施洗。} \BibleRef{玛 28:18}因此，异端者的虚假而邪恶的圣洗必须被我们拒绝，并以全然的憎恶加以驳斥，他们的口里说出的是毒药，而非生命，非属天的恩宠，而是对天主圣三的亵渎。因此显然，来到教会的异端者应当以健全而天主教的圣洗受洗，以便从他们僭越的亵渎中被净化，并借圣神的恩宠得以更新。\par

\Place{瓦加}的\Person{利博苏斯}说：\WorkTitle{福音}中主说：\BibleQuote{我是真理。} \BibleRef{若 14:6} 他没有说：\Quote{我是惯例。} 因此，真理既已显明，让惯例屈服于真理；如此，虽然过去任何人在教会内未曾习惯于为异端者授洗，现在就让他开始为他们授洗。\par

\Place{特贝斯德}的\Person{路基乌斯}说：我断定，那些亵渎、不义的异端者，以各种言语撕裂圣经圣洁可敬的话语，应受诅咒，因此必须为他们驱魔并授洗。\par

\Place{阿默德拉}的\Person{欧根尼乌斯}说：我也作同样的断定——异端者必须受洗。\par

还有一位\Place{阿马科拉}的\Person{斐理克斯}说：我自己也遵照圣经的权威，断定异端者必须受洗；而且，那些自认为在分裂者中已受洗的人，也必须受洗。因为，按照基督的警告，我们的圣洗之泉是私有的，让我们教会的一切敌人明白，它不能为别人所用。也不能由那唯一羊群的牧者将救恩之水赐给两个民族。因此，显然异端者和分裂者都不能领受任何属天的事物，因为他们竟敢从有罪的人、从教会之外的人那里领受。既然施予者没有资格，领受者当然也无益处。\par

还有一位\Place{穆祖利}的\Person{雅努阿留斯}说：我感到惊讶，既然所有人都承认只有一个圣洗，却没有人认识到这同一个圣洗的统一性。因为教会和异端是两回事，且是不同的事。如果异端者拥有圣洗，我们就未拥有；但如果我们拥有，异端者就不能拥有。但毫无疑问，唯独教会拥有基督的圣洗，因为唯独她拥有基督的恩宠和真理。\par

\Place{塔斯瓦尔特德}的\Person{阿德尔斐乌斯}说：有些人毫无理由地用虚假和嫉妒的言语攻击真理，说我们重洗，而教会并非为异端者重洗，而是为他们授洗。\par

\Place{莱普提米努斯}的\Person{德默特流斯}说：我们坚持一个圣洗，因为我们只为天主教会要求她自己的财产。但那些说异端者真实合法地施洗的人，正是他们自己造出了许多圣洗，而非一个。因为既然异端有许多，圣洗的数目也将按它们的数目计算。\par

\Place{提巴里斯}的\Person{味增爵}说：我们知道异端者比外邦人更恶劣。因此，如果他们在悔改后愿意来归向主，我们确实有真理的准则，就是主以神圣诫命命令他宗徒所说的：\Quote{去吧，因我的名按手，驱逐魔鬼。} 又在另一处说：\BibleQuote{去吧，教导万民，因圣父、圣子、圣神之名给他们授洗。} \BibleRef{玛 28:19} 因此，首先藉驱魔中按手，其次藉重生的圣洗，他们才能来到基督的应许中。否则，我认为不应如此行。\par

\Place{马克塔里斯}的\Person{马尔谷}说：异端者、真理的敌人和攻击者，将别人权力和屈尊之事归于自己，这不足为奇。但令我们惊奇的是，我们当中有些人，即真理的违背者，支持异端者，并与基督徒对立。因此我们判定，异端者必须受洗。\par

\Place{西西利巴}的\Person{萨提乌斯}说：如果异端者在圣洗中罪得赦免，他们就没有理由来到教会。因为在审判之日，正是罪受到惩罚，如果异端者已经获得罪的赦免，他们就不必惧怕基督的审判了。\par

\Place{戈尔}的\Person{维克多}说：既然只有在教会的圣洗中罪才得赦免，那么未经圣洗就接纳异端者共融的人，做了两件不合理的事：他既没有洁净异端者，又玷污了基督徒。\par

\Place{乌提卡}的\Person{奥勒留}说：既然宗徒说我们不可在别人的罪上有份，那么，与未经教会圣洗的异端者共融的人，岂非正是在别人的罪上有份？因此我断定，异端者必须受洗，好使他们获得罪的赦免，然后才能与他们共融。\par

\Place{日耳曼尼基雅纳}的\Person{扬布斯}说：那些赞同异端者圣洗的人，就是否定我们的圣洗，因为他们否认那些在教会之外（我不说被洗净，而是被玷污）的人，应当在教会内受洗。\par

\Place{鲁库玛}的\Person{路基雅努斯}说：经上记载：\BibleQuote{天主见光好，就将光与黑暗分开。} \BibleRef{创 1:4} 如果光与黑暗可以相和，那么我与异端者之间也可有共同之处。因此我断定，异端者必须受洗。\par

\Place{卢佩尔基雅纳}的\Person{贝拉基雅努斯}说：经上记载：\BibleQuote{要么主是天主，要么巴耳是天主。} \BibleRef{列上 18:21} 因此，在目前的情况下，要么教会是教会，要么异端是教会。另一方面，如果异端不是教会，教会的圣洗怎能存在于异端者中？\par

\Place{米迪拉}的\Person{雅德尔}说：我们知道在天主教会中只有一个圣洗，因此我们不应接纳异端者，除非他在我们中间受洗，以免他认为自己在天主教会之外已经受洗。\par

还有一位\Place{马腊扎纳}的\Person{斐理克斯}说：只有一个信德，一个圣洗，但属于天主教会，唯独她有权施洗。\par

\Place{奥巴}的\Person{保禄}说：如果有人不坚持教会的信德和真理，我并不困扰，因为宗徒说：\BibleQuote{即使有些人不信，难道他们的不信会使天主的信实失效吗？断乎不能！天主是真实的，人都是虚谎的。} \BibleRef{罗 3:3} 但如果天主是真实的，圣洗的真理怎能存在于异端者中，在天主不在的地方？\par

\Place{狄奥尼西雅纳}的\Person{蓬波纽斯}说：显然，异端者不能施洗并赦免罪过，因为他们在地上没有能力束缚或释放任何事物。\par

\Person{韦南修} 说：若一位丈夫远赴异乡，将自己的妻子托付给朋友监护，那位朋友必将尽心竭力守护受托付的妻子，免得她的贞洁与圣德被任何人玷污。主基督、我们的天主，往父那里去了，将祂的新娘托付给我们。我们该守护她，使她无玷无瑕，还是将她的完整与贞洁出卖给奸夫和败坏者？因为那将教会的圣洗与异端者共用的，就是把基督的净配出卖给奸夫。\par

\Person{阿希姆努斯} 说：我们领受了同一的圣洗，我们也持守并实践这同一的圣洗。但谁若说异端者也能合法地施洗，就是制造了两个圣洗。\par

\Person{撒杜尼努斯} 说：若异端者可以施洗，那么行不法之事的人就被谅解和辩护了；我也不明了为何基督称他们为仇敌，或宗徒称他们为假基督。\par

\Person{撒杜尼努斯} 说：外邦人虽然崇拜偶像，却仍然认识并承认有一位至高的天主为圣父和创造者。\Person{马尔西昂} 亵渎祂，而有些人竟不羞于认可马尔西昂的圣洗。这样的司铎，既不洗天主的仇敌，又与他们共融，他们是如何遵守或维护天主的司祭职呢？\par

\Person{马尔塞鲁} 说：既然罪过只在教会的圣洗中才得赦免，那么不洗异端者的人，便是与罪人共融。\par

\Person{依肋乃} 说：若教会不洗异端者，理由是他已经受洗了，那便是更大的异端。\par

\Person{杜纳图} 说：我知道只有一个教会及其一个圣洗。若有人说圣洗的恩宠在异端者那里，他必须先表明并证明教会在他们当中。\par

\Person{佐西默} 说：当真理得到启示时，容错误让位于真理；因为\Person{伯多禄} 先前施行割损，当\Person{保禄} 宣讲真理时，他也顺从了保禄。\par

\Person{犹利安} 说：经上记载：\Quote{人不能领受什么，除非从上头赐给了他。} 若异端来自上天，它也能给人圣洗。\par

\Person{福斯托} 说：让那些祖护异端者的人不要自欺。谁若为异端者干涉教会的圣洗，便是使他们成为基督徒，而使我们成为异端者。\par

\Person{革米尼} 说：我们的某些同僚或许将异端者看得比他们自己还重，但他们不能这样看待我们；因此我们一旦决定了的事，我们就持守——即我们为那些从异端者那里来的人施洗。\par

\Person{罗加提亚努斯} 说：基督建立了教会；魔鬼则建立了异端。撒旦的会堂怎能拥有基督的圣洗？\par

\Person{特拉皮乌} 说：谁将教会的圣洗让与并出卖给异端者，他对于基督的净配除了是\Person{犹达斯} 之外还是什么？\par

另有\Person{路爵} 说：经上记载：\Quote{天主不听罪人。} 一个身为罪人的异端者，在圣洗中怎能被垂听？\par

另有\Person{斐理斯} 说：在未经教会圣洗而接纳异端者的事上，愿无人将习惯置于理性和真理之上，因为理性和真理总是排除习惯。\par

\Person{撒杜尼努斯} 说：若假基督能给人基督的恩宠，异端者也能施洗，因为他们被称为假基督。\par

\Person{君托} 说：拥有某物的人才能给予某物。但异端者显然一无所有，他们能给予什么呢？\par

\Person{犹利安} 说：若一个人能服侍两个主人，天主和玛门，那么圣洗也能服侍两个主人，基督徒和异端者。\par

\Person{特纳克斯} 说：圣洗只有一个，但它是教会的。哪里没有教会，哪里就没有圣洗。\par

\Person{维克多} 说：经上记载：\Quote{天主是一位，基督是一位，教会是一个，圣洗是一个。} 那么，在没有天主、基督和唯一教会的地方，人怎能受洗呢？\par

\Person{杜纳突罗} 说：我也一直这样认为，即异端者在教会之外不能获得任何东西，当他们归向教会时，必须受洗。\par

\Person{维鲁罗} 说：身为异端者的人不能给予他没有的东西；分裂者更是如此，因为他已失去了从前拥有的。\par

\Person{普登提亚努斯} 说：亲爱的弟兄们，我主教职位的初任使我等待长老们的判断。因为显然异端没有任何东西，也不能有任何东西。因此，若有人从他们那里来，最公义的规定是他们必须受洗。\par

\Person{伯多禄} 说：既然在天主教会中只有一个圣洗，显然人不能在教会之外受洗。因此我判定，那些在异端或分裂中受浸的人，当他们来到教会时，应当受洗。\par

另有\Person{路爵} 说：按照我的心意和圣神的引导，既然有一位天主，我们主耶稣基督的圣父，一位基督，一个希望，一位圣神，一个教会，那么也应当只有一个圣洗。因此我说，凡异端者所发起或成就的任何事，都应当废止，那些从那里来的人必须在教会中受洗。\par

\Person{斐理斯} 说：我判定，按照圣经的训示，凡在教会之外由异端者非法受洗的人，当他愿意投靠教会时，应当在圣洗依法施予的地方获得圣洗的恩宠。\par

\Person{普西罗} 说：我相信除非在天主教会中，没有使人得救的圣洗。凡在天主教会之外的，都是虚假。\par

\Person{萨尔维亚努斯} 说：异端者确实一无所有，因此他们来到我们这里，为要获得他们所没有的。\par

\Person{霍诺拉图} 说：既然基督是真理，我们宁愿跟随真理而不跟随习惯；因此我们应当用教会的圣洗圣化异端者，因为他们来到我们这里，是因为他们在外面什么也不能领受。\par

\Person{维克多} 说：正如你们自己也知道的，我被任命为主教不久，因此我等候了前任的决定。因此我的看法是，凡从异端来的人，毫无疑问应当受洗。\par

\Signature{马苏拉的克拉鲁斯}说：主耶稣基督的声明是清楚的，当祂派遣祂的宗徒时，将祂由圣父所获的权能唯独赐予了他们；我们继承了他们的位置，以同一权能治理主的教会，并为信徒的信仰施洗。因此，异端者既在外面没有权能，又没有基督的教会，不能以祂的洗礼为任何人施洗。\par

\Signature{塔姆贝的塞孔迪亚努斯}说：我们不应以我们的冒失欺骗异端者，以致那些未在我们主耶稣基督的教会中受洗、并因此未获得罪赦的人，在审判之日来临时，将归咎于我们，说因着我们他们未受洗，未获得天主恩宠的赦免。因此，既然只有一个教会和一个洗礼，当他们归向我们时，应当连同教会一起获得教会的洗礼。\par

\Signature{楚拉比的奥勒留}说：宗徒若望在他的书信中定下法则说：\BibleQuote{若有人来到你们那里，没有基督的道理，不要接他到家里，也不要向他说：安好。因为向他说安好的，就与他的恶行有分了。} 怎能如此轻率地接纳这样的人进入天主的家呢？他们连进入我们的私宅都被禁止。或者，我们怎能不与教会的洗礼就与他们共融呢？我们若是只向他们说声安好，就与他们的恶行有分了。\par

\Signature{杰梅利的利特乌斯}说：如果瞎子领瞎子，两人都掉在坑里。既然很明显，异端者不能给任何人光明，因为他们自己是瞎子，他们的洗礼就没有效力。\par

\Signature{奥埃的纳塔利斯}说：我本人，以及萨布拉塔的庞培，还有莱普提斯马格纳的迪奥伽——他虽身不在场，但心神同在，并指示我——都判断与我们的同僚一样，即异端者除非领受教会的洗礼，否则不能与我们共融。\par

\Signature{纳波利斯的尤尼乌斯}说：我不收回我们曾决定的判决，即应给来到教会的异端者施洗。\par

\Signature{迦太基的西彼廉}说：写给我们的同僚犹巴亚努斯的信充分表达了我的意见，即根据福音和宗徒的见证，被称为基督的敌人和\OriginalTerm{敌基督者}{Antichrists}的异端者，当他们来到教会时，必须领受教会唯一的洗礼，使他们从敌人变为朋友，从敌基督者变为基督徒。\par

\SourceRecord{Translated by Robert Ernest Wallis.；Ante-Nicene Fathers；Vol. 5.；Edited by Alexander Roberts, James Donaldson, and A. Cleveland Coxe.；Buffalo, NY: Christian Literature Publishing Co.,；1886.；<http://www.newadvent.org/fathers/0508.htm>.}
